\documentclass[10pt,a4paper]{amsart}
\usepackage{amsmath,amssymb,amsthm,graphicx}
% Retained minimal preamble for the Stacks Project chapter flat.tex (More on Flatness).
% Verbatim excerpt of the chapter preamble: Xy-pic setup (native preamble lines 14-18) and the native
% theorem declarations with the mathematical macros (native preamble lines 154-217). Copied mechanically;
% nothing added, removed or redefined. The wrapper class and the amsmath/amssymb/amsthm/graphicx base
% packages are supplied by the assembler, so the class/externaldocument part of the native preamble is omitted.
% For commutative diagrams we use Xy-pic
\usepackage[all]{xy}

% We use 2cell for 2-commutative diagrams.
\xyoption{2cell}
% Theorem environments.
%
\theoremstyle{plain}
\newtheorem{theorem}[subsection]{Theorem}
\newtheorem{proposition}[subsection]{Proposition}
\newtheorem{lemma}[subsection]{Lemma}

\theoremstyle{definition}
\newtheorem{definition}[subsection]{Definition}
\newtheorem{example}[subsection]{Example}
\newtheorem{exercise}[subsection]{Exercise}
\newtheorem{situation}[subsection]{Situation}

\theoremstyle{remark}
\newtheorem{remark}[subsection]{Remark}
\newtheorem{remarks}[subsection]{Remarks}

\numberwithin{equation}{subsection}

% Macros
%
\def\lim{\mathop{\mathrm{lim}}\nolimits}
\def\colim{\mathop{\mathrm{colim}}\nolimits}
\def\Spec{\mathop{\mathrm{Spec}}}
\def\Hom{\mathop{\mathrm{Hom}}\nolimits}
\def\Ext{\mathop{\mathrm{Ext}}\nolimits}
\def\SheafHom{\mathop{\mathcal{H}\!\mathit{om}}\nolimits}
\def\SheafExt{\mathop{\mathcal{E}\!\mathit{xt}}\nolimits}
\def\Sch{\mathit{Sch}}
\def\Mor{\mathop{\mathrm{Mor}}\nolimits}
\def\Ob{\mathop{\mathrm{Ob}}\nolimits}
\def\Sh{\mathop{\mathit{Sh}}\nolimits}
\def\NL{\mathop{N\!L}\nolimits}
\def\CH{\mathop{\mathrm{CH}}\nolimits}
\def\proetale{{pro\text{-}\acute{e}tale}}
\def\etale{{\acute{e}tale}}
\def\QCoh{\mathit{QCoh}}
\def\Ker{\mathop{\mathrm{Ker}}}
\def\Im{\mathop{\mathrm{Im}}}
\def\Coker{\mathop{\mathrm{Coker}}}
\def\Coim{\mathop{\mathrm{Coim}}}

% Boxtimes
%
\DeclareMathSymbol{\boxtimes}{\mathbin}{AMSa}{"02}

%
% Macros for moduli stacks/spaces
%
\def\QCohstack{\mathcal{QC}\!\mathit{oh}}
\def\Cohstack{\mathcal{C}\!\mathit{oh}}
\def\Spacesstack{\mathcal{S}\!\mathit{paces}}
\def\Quotfunctor{\mathrm{Quot}}
\def\Hilbfunctor{\mathrm{Hilb}}
\def\Curvesstack{\mathcal{C}\!\mathit{urves}}
\def\Polarizedstack{\mathcal{P}\!\mathit{olarized}}
\def\Complexesstack{\mathcal{C}\!\mathit{omplexes}}
% \Pic is the operator that assigns to X its picard group, usage \Pic(X)
% \Picardstack_{X/B} denotes the Picard stack of X over B
% \Picardfunctor_{X/B} denotes the Picard functor of X over B
\def\Pic{\mathop{\mathrm{Pic}}\nolimits}
\def\Picardstack{\mathcal{P}\!\mathit{ic}}
\def\Picardfunctor{\mathrm{Pic}}
\def\Deformationcategory{\mathcal{D}\!\mathit{ef}}
\usepackage{cleveref}
\crefname{originallemma}{original lemma}{original lemmas}
\Crefname{originallemma}{Original lemma}{Original lemmas}
\makeatletter
\newlabel{lemma-lift-etale}{{057R}{source}{}{sourcepointer-1}{}}
\newlabel{lemma-lift-etale@cref}{{[originallemma][0][]057R}{[1][0][]source}{}{}{}}
\makeatother
\crefname{originallemma}{original lemma}{original lemmas}
\Crefname{originallemma}{Original lemma}{Original lemmas}
\makeatletter
\newlabel{lemma-sheaf-lives-on-subscheme}{{057Q}{source}{}{sourcepointer-2}{}}
\newlabel{lemma-sheaf-lives-on-subscheme@cref}{{[originallemma][0][]057Q}{[1][0][]source}{}{}{}}
\makeatother
\crefname{originallemma}{original lemma}{original lemmas}
\Crefname{originallemma}{Original lemma}{Original lemmas}
\makeatletter
\newlabel{more-morphisms-lemma-local-local-structure-finite-type}{{057K}{source}{}{sourcepointer-3}{}}
\newlabel{more-morphisms-lemma-local-local-structure-finite-type@cref}{{[originallemma][0][]057K}{[1][0][]source}{}{}{}}
\makeatother
\crefname{originallemma}{original lemma}{original lemmas}
\Crefname{originallemma}{Original lemma}{Original lemmas}
\makeatletter
\newlabel{more-morphisms-lemma-local-local-structure-finite-type-affine}{{057L}{source}{}{sourcepointer-4}{}}
\newlabel{more-morphisms-lemma-local-local-structure-finite-type-affine@cref}{{[originallemma][0][]057L}{[1][0][]source}{}{}{}}
\makeatother
\crefname{originallemma}{original lemma}{original lemmas}
\Crefname{originallemma}{Original lemma}{Original lemmas}
\makeatletter
\newlabel{morphisms-lemma-etale-permanence}{{02GW}{source}{}{sourcepointer-5}{}}
\newlabel{morphisms-lemma-etale-permanence@cref}{{[originallemma][0][]02GW}{[1][0][]source}{}{}{}}
\makeatother
\title{Elementary etale localization producing a one step devissage of a finite type module}
\author{The Stacks Project Authors}
\date{stacks-project:flat.tex:local-snapshot-20261011; GNU FDL 1.2 or later}
\begin{document}
\maketitle
\noindent\textit{Editorial note.} All delivered mathematical text is the author's own, copied from the pinned Stacks Project chapter file flat.tex; no statement or proof step was rewritten, reordered or removed.\par\medskip
\noindent\textit{Editorial note.} The preamble is the chapter's own retained minimal preamble (sources/stacks-flat/retained-preamble.tex), a verbatim excerpt of the native Xy-pic setup and of the native theorem declarations and mathematical macros. The native class line, the xr-hyper externaldocument block and the reference/slogan/history display environments are omitted because the assembler supplies a plain amsart wrapper and the excerpt uses none of those environments; no macro was invented or redefined.\par\medskip
\noindent\textit{Editorial note.} Locators of the cross-chapter references point either at the \textbackslash{}label line of the referenced result inside the pinned chapter file (same-chapter references) or at a verbatim excerpt of the project's own tag list, sources/stacks-flat/referenced-tags.txt, for results living in other chapters, each recorded with its stable Stacks tag.\par\medskip
\noindent\textit{Editorial note.} The notion of an elementary etale neighbourhood of pointed schemes is defined in earlier chapters of the same project and not in this file, so it appears among the standard prerequisites rather than as a context span.\par\medskip
\noindent\textit{Allowed standard prerequisites (editorial).}\begin{itemize}
\item Elementary etale neighbourhoods of pointed schemes. The statement quantifies over elementary etale neighbourhoods (S', s') -> (S, s) and (X', x') -> (X, x). The construction of the diagram and of the closed immersion into it is the retained obligation; the notion itself is standard input.
\item Support of a finite type module as a closed subscheme cut out by its annihilator. Only the existence of Z -> X with i\_*G = F and Z\_s = Supp(F\_s) is taken from this pointer; the gluing of that closed immersion with the etale localization of the source stays inside the counted proof.
\item Local structure of a morphism locally of finite type over an affine base. Used only for the case X\_s = Supp(F\_s) at proof lines 534-549; the general case at proof lines 551-594 is retained in full and is the counted core.
\end{itemize}\medskip
\ifdefined\hypertarget\hypertarget{sourcepointer-1}{}\fi
\noindent\textit{Source reference.} Original lemma 057R: sources/stacks-flat/pinned-original-source.tex, source lines 84--84. Same-chapter navigation pointer to the etale lifting lemma invoked at proof lines 566-579; the counted construction (case split, support closed immersion, gluing of the diagram) is retained verbatim in the proof spans, so this pointer replaces nothing.\par\medskip
\ifdefined\hypertarget\hypertarget{sourcepointer-2}{}\fi
\noindent\textit{Source reference.} Original lemma 057Q: sources/stacks-flat/pinned-original-source.tex, source lines 446--446. Same-chapter pointer for the existence of the closed immersion cutting out the support, used at proof line 554 as an input; the retained proof still carries the case analysis and the etale-lifting step that constitute the counted obligation.\par\medskip
\ifdefined\hypertarget\hypertarget{sourcepointer-3}{}\fi
\noindent\textit{Source reference.} Original lemma 057K: sources/stacks-flat/referenced-tags.txt, source lines 8--8. Cross-chapter standard prerequisite (local structure of a finite type morphism) invoked at proof lines 536-539 for the case X\_s = Supp(F\_s); the general case of the counted proof is retained in full.\par\medskip
\ifdefined\hypertarget\hypertarget{sourcepointer-4}{}\fi
\noindent\textit{Source reference.} Original lemma 057L: sources/stacks-flat/referenced-tags.txt, source lines 9--9. Cross-chapter standard prerequisite invoked together with 057K at proof lines 536-539; it supplies the affine local model only, not the counted gluing argument.\par\medskip
\ifdefined\hypertarget\hypertarget{sourcepointer-5}{}\fi
\noindent\textit{Source reference.} Original lemma 02GW: sources/stacks-flat/referenced-tags.txt, source lines 4--4. Cross-chapter notational pointer recording that the first morphism Z' -> Z x\_S S' in the displayed composite is etale (proof line 575-576); the surrounding argument is retained.\par\medskip
\begin{lemma}
\label{lemma-elementary-devissage}
Let $f : X \to S$ be morphism of schemes which is locally of finite type.
Let $\mathcal{F}$ be a finite type quasi-coherent $\mathcal{O}_X$-module.
Let $x \in X$ with image $s = f(x)$ in $S$.
Set $\mathcal{F}_s = \mathcal{F}|_{X_s}$ and
$n = \dim_x(\text{Supp}(\mathcal{F}_s))$.
Then we can construct
\begin{enumerate}
\item elementary \'etale neighbourhoods $g : (X', x') \to (X, x)$,
$e : (S', s') \to (S, s)$,
\item a commutative diagram
$$
\xymatrix{
X \ar[dd]_f & X' \ar[dd] \ar[l]^g & Z' \ar[l]^i \ar[d]^\pi \\
& & Y' \ar[d]^h \\
S & S' \ar[l]_e & S' \ar@{=}[l]
}
$$
\item a point $z' \in Z'$ with $i(z') = x'$, $y' = \pi(z')$, $h(y') = s'$,
\item a finite type quasi-coherent $\mathcal{O}_{Z'}$-module $\mathcal{G}$,
\end{enumerate}
such that the following properties hold
\begin{enumerate}
\item $X'$, $Z'$, $Y'$, $S'$ are affine schemes,
\item $i$ is a closed immersion of finite presentation,
\item $i_*(\mathcal{G}) \cong g^*\mathcal{F}$,
\item $\pi$ is finite and $\pi^{-1}(\{y'\}) = \{z'\}$,
\item the extension $\kappa(y')/\kappa(s')$ is purely transcendental,
\item $h$ is smooth of relative dimension $n$
with geometrically integral fibres.
\end{enumerate}
\end{lemma}

\begin{proof}
Let $V \subset S$ be an affine neighbourhood of $s$.
Let $U \subset f^{-1}(V)$ be an affine neighbourhood of $x$.
Then it suffices to prove the lemma for $f|_U : U \to V$ and
$\mathcal{F}|_U$. Hence in the rest of the proof we assume that
$X$ and $S$ are affine.

\medskip\noindent
First, suppose that $X_s = \text{Supp}(\mathcal{F}_s)$, in particular
$n = \dim_x(X_s)$. Apply
More on Morphisms,
Lemmas \ref{more-morphisms-lemma-local-local-structure-finite-type} and
\ref{more-morphisms-lemma-local-local-structure-finite-type-affine}.
This gives us a commutative diagram
$$
\xymatrix{
X \ar[dd] & X' \ar[l]^g \ar[d]^\pi \\
& Y' \ar[d]^h  \\
S & S' \ar[l]_e
}
$$
and point $x' \in X'$. We set $Z' = X'$, $i = \text{id}$, and
$\mathcal{G} = g^*\mathcal{F}$ to obtain a solution in this case.

\medskip\noindent
In general choose a closed immersion $Z \to X$ and a sheaf
$\mathcal{G}$ on $Z$ as in
Lemma \ref{lemma-sheaf-lives-on-subscheme}.
Applying the result of the previous paragraph to $Z \to S$ and
$\mathcal{G}$ we obtain a diagram
$$
\xymatrix{
X \ar[dd]_f & Z \ar[l] \ar[dd]_{f|_Z} & Z' \ar[l]^g \ar[d]^\pi \\
& & Y' \ar[d]^h \\
S & S \ar@{=}[l] & S' \ar[l]_e
}
$$
and point $z' \in Z'$ satisfying all the required properties.
We will use
Lemma \ref{lemma-lift-etale}
to embed $Z'$ into a scheme \'etale over $X$. We cannot apply the lemma directly
as we want $X'$ to be a scheme over $S'$. Instead we
consider the morphisms
$$
\xymatrix{
Z' \ar[r] & Z \times_S S' \ar[r] & X \times_S S'
}
$$
The first morphism is \'etale by
Morphisms, Lemma \ref{morphisms-lemma-etale-permanence}.
The second is a closed immersion as a base change of a closed immersion.
Finally, as $X$, $S$, $S'$, $Z$, $Z'$ are all affine we may apply
Lemma \ref{lemma-lift-etale}
to get an \'etale morphism of affine schemes $X' \to X \times_S S'$ such that
$$
Z' = (Z \times_S S') \times_{(X \times_S S')} X' = Z \times_X X'.
$$
As $Z \to X$ is a closed immersion of finite presentation, so is $Z' \to X'$.
Let $x' \in X'$ be the point corresponding to $z' \in Z'$.
Then the completed diagram
$$
\xymatrix{
X \ar[dd] & X' \ar[dd] \ar[l] & Z' \ar[l]^i \ar[d]^\pi \\
& & Y' \ar[d]^h \\
S & S' \ar[l]_e & S' \ar@{=}[l]
}
$$
is a solution of the original problem.
\end{proof}

\end{document}
